\documentclass[11pt]{article}
\usepackage[a4paper,margin=27mm]{geometry}
\usepackage{amsmath,amssymb,amsthm,hyperref}
\newtheorem{theorem}{Theorem}
\newtheorem{lemma}[theorem]{Lemma}
\newtheorem{proposition}[theorem]{Proposition}
\newcommand{\Aut}{\operatorname{Aut}}
\newcommand{\Fix}{\operatorname{Fix}}
\title{Counting maximal-subgroup pairs from fixed points}
\author{Working proof for the remaining claims of L036}
\date{}
\begin{document}\maketitle
This note supplies reductions for the remaining claims of Leemans and
Rodrigues, \emph{On primitive geometries of rank two}, Art Discrete Appl.
Math.\ 2 (2019), P1.08. It contains no accepted numerical application yet.
The source's equivalence is the action of $\Aut(G)$ on unordered pairs of
distinct maximal subgroups of $G$. A pair determines a firm, connected,
flag-transitive primitive coset geometry. Conversely, the source's
geometries have precisely this form. These facts and the Rudvalis lower
bound are already proved in the accepted principal correction.

\section{The conjugation action, including outer automorphisms}
For a finite group $A$, choose representatives $H_1,\ldots,H_r$ of all
conjugacy classes of subgroups. Write
\[
 m_{ij}=|(A/H_i)^{H_j}|.
\]
Thus the rows of the table of marks give fixed points on transitive
$A$-sets. In particular, $m_{ij}>0$ if and only if $H_j$ is contained
in an $A$-conjugate of $H_i$. Also
\[
 m_{i1}=[A:H_i],\qquad m_{ii}=[N_A(H_i):H_i],
\]
where $H_1=1$. The first equality follows by taking fixed points of the
trivial group. For the second, a coset fixed by $H_i$ corresponds to a
conjugate of $H_i$ containing $H_i$; equal orders force equality, and
the set of such cosets is $N_A(H_i)/H_i$.

\begin{lemma}\label{normalizer}
Let $G$ be a finite nonabelian simple group and let
$G\leq A\leq\Aut(G)$, where $G$ is identified with its inner
 automorphism group. Suppose $[A:G]\leq2$. If $M$ is maximal in $G$,
then $N_A(M)$ is determined, up to $A$-conjugacy, by the orders and
containment relation in the table of marks of $A$.
\end{lemma}
\begin{proof}
The group $M$ is nontrivial: a nonabelian simple group cannot have the
trivial group as a maximal subgroup. Maximality gives $N_G(M)=M$ or
$G$. The latter would make $M$ a nontrivial proper normal subgroup,
so $N_G(M)=M$. The map $N_A(M)\to A/G$ has kernel $M$. Consequently
$[N_A(M):M]$ is one or two. This index is the diagonal mark at $M$.
If it is one, the normalizer is $M$.

Suppose it is two. Any subgroup $K$ of order $2|M|$ containing $M$
normalizes $M$, because a subgroup of index two is normal. Hence
$K\leq N_A(M)$, and equal orders give $K=N_A(M)$. At least one such
subgroup exists, namely the normalizer itself. Thus precisely one
conjugacy class of subgroups of order $2|M|$ can contain a conjugate
of $M$. Both its existence and its class are read from the table.
\end{proof}

If $[A:G]=2$, the class of $G$ is the unique subgroup class of order
$|A|/2$. Indeed, any subgroup $K$ of that index is normal in $A$.
The intersection $G\cap K$ is normal in $G$ and has index at most two
there. Simplicity and $|G|>2$ imply $G\leq K$, and then $G=K$.
The maximal subgroups of $G$, up to $A$-conjugacy, are exactly the
maximal proper classes below this class in the containment relation.
Because $G$ is normal, containment in a conjugate of $G$ is simply
containment in $G$. If $A=G$, use the final class instead.

\begin{theorem}\label{count}
Under the hypotheses of Lemma~\ref{normalizer}, let $\mathcal M$
contain one representative of every $A$-conjugacy class of maximal
subgroups of $G$. For a conjugacy class representative $a\in A$, put
\[
 \psi(a)=\sum_{M\in\mathcal M}m_{N_A(M),\langle a\rangle},
\]
where the subscripts denote the corresponding subgroup classes.
Then the number of $A$-orbits on unordered pairs of distinct maximal
subgroups of $G$ is
\[
 \frac12\sum_{[a]_A}
 \frac{\psi(a)^2-2\psi(a)+\psi(a^2)}{|C_A(a)|}.                 \tag{1}
\]
Thus a complete table of marks, conjugacy-class centralizers, a square
map, and the fusion taking an element class to its cyclic subgroup
class determine this number by exact integer arithmetic.
\end{theorem}
\begin{proof}
The conjugation orbit of $M$ is the transitive $A$-set
$A/N_A(M)$. Its points fixed by $a$ are its points fixed by
$\langle a\rangle$, hence are counted by the displayed mark.
The disjoint union of these orbits is the set of all maximal subgroups
of $G$. Therefore $\psi(a)$ is its fixed-point count.

An unordered pair fixed by $a$ is either a pair of fixed points or
a two-cycle. There are $\binom{\psi(a)}2$ of the first kind and
$(\psi(a^2)-\psi(a))/2$ of the second. Burnside's orbit-counting
lemma, grouped by element conjugacy classes, gives (1), since the
class of $a$ has size $|A|/|C_A(a)|$.
\end{proof}
The formula uses the full action of $A$. Replacing it by the action of
$G$ when outer automorphisms are present does not compute the source's
isomorphism classes.

When $A=G$, a smaller character-table certificate suffices. Since
$N_G(M)=M$, the required action is the union of the coset actions
$G/M$. If the classes of $M$ fuse to those of $G$, then
\[
 |(G/M)^g|=|C_G(g)|
     \sum_{[h]_M\,\mapsto\,[g]_G}\frac{1}{|C_M(h)|}.          \tag{2}
\]
To prove (2), count the elements $x\in G$ with $x^{-1}gx\in M$.
Each suitable element of $M$ has $|C_G(g)|$ preimages. Dividing by
$|M|$ counts cosets; expanding the suitable elements into $M$-classes
gives the stated expression.

\section{Suborbit lengths from products of marks}
\begin{theorem}\label{product}
Let $H,K\leq G$ and let $m_L$ denote the mark row of $G/L$.
There are unique nonnegative integers $c_L$, indexed by the conjugacy
classes of subgroups of $G$, such that
\[
              m_H\,m_K=\sum_{[L]_G}c_Lm_L,                 \tag{3}
\]
where multiplication on the left is coordinatewise. The multiset of
orbit lengths of $H$ on $G/K$ consists of $c_L$ copies of $[H:L]$
for each $L$ with $c_L>0$. In particular, $|L|$ divides $|H|$ for
every such coefficient.
\end{theorem}
\begin{proof}
Fixed points commute with Cartesian products. Thus the left side of
(3) is the fixed-point vector of the diagonal action of $G$ on
$G/H\times G/K$. Decompose this finite $G$-set into its transitive
orbits. Each orbit is isomorphic to $G/L$ for the stabilizer of one
of its points; collecting conjugate stabilizers gives (3) with
nonnegative integral coefficients.

Order the subgroup representatives by increasing order, with an
arbitrary order among equal orders. Containment makes the table of
marks triangular, with nonzero diagonal. Its rows are linearly
independent over $\mathbb Q$, proving uniqueness and giving an exact
triangular solution for (3).

Every diagonal orbit meets $\{H\}\times G/K$. Its intersection
with this fiber is one $H$-orbit: if a group element sends a point of
the fiber to another, it fixes the first coordinate and belongs to
$H$. The stabilizer of $(H,gK)$ is $H\cap gKg^{-1}$. Hence the
corresponding $H$-orbit has length $|H|/|H\cap gKg^{-1}|$.
Its stabilizer is conjugate to the $L$ attached to the diagonal
orbit, proving the formula and divisibility assertion.
\end{proof}
This theorem recovers the full multiset of suborbit lengths, including
singletons. For two conjugacy classes of maximal subgroups, singleton
orbits occur precisely when the two actual subgroups coincide. The
suborbit multiset alone does not identify which orbits are interchanged
by duality or by outer automorphisms; Theorem~\ref{count} handles that
separate equivalence. Nor does it recover diameters, girths or design
parameters of the individual incidence graphs.

\section{Attribution and remaining certificate obligations}
Tables of marks, their injective fixed-point map, the product (3),
induction of trivial characters, and Burnside's lemma are classical.
Our proposed contribution is their application to the source's complete
list of counting and suborbit claims, with an explicit normalizer step
for the outer automorphisms. The GAP Table of Marks Library supplies a
complete subgroup table as an attributed premise; internal consistency
checks do not prove that the table is complete or belongs to the named
group. The Character Table Library supplies the named character tables,
complete maximal-subgroup list and compatible fusions when (2) is used.
The installed documentation explicitly states this meaning of
\texttt{Maxes}. All library editions and selected files must be bound
before a numerical application is accepted.

The accepted $M_{22}$ census in the existing L073 recovery already gives
81 pair classes under the required full automorphism equivalence. It
may be cited without redoing its enumeration. The other retained raw
character totals have not yet been accepted. The source's 37 $M_{11}$
incidence graphs and their design parameters still require actual
representatives and finite checks. Reported historical timings are
observations of the original implementation, not mathematical claims
recovered by the fixed-point formulas.
\end{document}
